\documentclass[10pt,a4paper]{amsart}
\usepackage[margin=19mm]{geometry}
\usepackage{amsmath,amssymb,mathtools}
\usepackage[scr=rsfs,cal=euler]{mathalfa}
\usepackage{xfrac}
\usepackage[unicode,hidelinks,bookmarks=false]{hyperref}
\newcommand{\ie}{i.e.\ }
\newcommand{\cf}{cf.\ }
\newcommand{\resp}{respectively\ }
\newcommand{\Subsection}{\subsection}
\providecommand{\nobreakdash}{\nobreak\hbox{-}}
\setlength{\emergencystretch}{2em}
\multlinegap0pt
\usepackage{amssymb}
\usepackage{tikz}
\newtheorem{lemma}{Lemma}
\newcommand{\M}  {\mathcal{M}}
\newcommand{\N}  {{\mathbb N}}
\newcommand{\R}  {{\mathbb R}}
\newcommand{\U}  {\mathcal{U}}
\newcommand{\W}  {\mathcal{W}}
\newcommand{\bg}{{\boldsymbol{\gamma}}}
\title{Embeddings of Reproducing Kernel Hilbert Spaces with General Weights}
\author{Michael Gnewuch, Peter Kritzer, Klaus Ritter}
\date{Source: arXiv:2604.27160v1 (CC BY 4.0)}
\begin{document}
\maketitle

\noindent\emph{Source and licence.} Embeddings of Reproducing Kernel Hilbert Spaces with General Weights. Version arXiv:2604.27160v1 (CC BY 4.0), distributed by arXiv under the Creative Commons Attribution 4.0 International licence, http://creativecommons.org/licenses/by/4.0/, as recorded in the arXiv licence metadata for this exact version and shown on the abstract page https://arxiv.org/abs/2604.27160v1. Full licence notice: \texttt{licenses/arxiv-2604.27160-CC-BY-4.0.txt}.\par\smallskip

\noindent\textit{Editorial scope.} Counted unit: the article's lemma Lemma33 (source0.tex lines 3234--3242). The article's complete original proof of it is delivered with it (source0.tex lines 3244--3318, one \texttt{proof} environment of 75 lines). No mathematics is written, repaired or re-derived: the statement and the proof are the article's own lines, in the article's own notation, and no retained line is substituted or re-flowed.  Retained uncounted as the source-local context the proof needs: lines 3166--3174.  Nothing was omitted from any retained span, so the package carries no omission record.  No retained line contains a citation, so the wrapper carries no bibliography and no bibliography entry.  No retained line contains a figure environment, an image-inclusion command or drawing code, so no image is delivered and the package declares no asset dependency.  Editorial formatting only, in addition: an AMS \texttt{amsart} preamble that restates the article's own declarations and macros used by the retained lines, namely the environments \texttt{lemma} on separate counters, together with the standard packages those lines need.  The retained environments print in this document as Lemma 1 in order of appearance, whereas the article prints the same items with the numbering of its own full document.  The complete native source of the article as distributed in the arXiv e-print for this exact version is bundled unchanged as \texttt{sources/199449/source0.tex}.
\begin{lemma}\label{Lemma32}
Let $d\in\N$. A function $f \colon [0,1]^d \to \R$ is 
completely monotonically decreasing if and only if 
\[
\sum_{w \subseteq v} (-1)^{|w|} f(x_{-w}\!:\!y_w)\ge 0
\]
for all non-empty $v\in \U_d$ and 
$x, y\in [0,1]^d$ such that $x_j \leq y_j$ for every $j \in v$.
\end{lemma}

\begin{lemma}\label{Lemma33}
Let $d\in \N$.
The mapping $\Phi$ defines a bijection from $\M_d$ onto 
\[
\M_d^\prime := \{ f|_{\{0,1\}^d} : 
\text{$f \colon [0,1]^d \to {[0,\infty)}$ completely monotonically
decreasing}\}. 
\]
\end{lemma}

\begin{proof}
For $f \colon [0,1]^d \to {[0,\infty)}$ let $\bg \in \W_d$ be
defined by 
\begin{equation}\label{g88}
\Phi \bg  = f|_{\{0,1\}^d}. 
\end{equation}
Furthermore, let $u,v \subseteq [d]$ with $u \cap v = \emptyset$ and
let $x,y \in \{0,1\}^d$ be defined by
$\varphi(x) = u$ and $\varphi(y) = v$.
Note that $u \cup w = \varphi(x_{-w}\!:\!y_w)$ for every $w \subseteq v$. 
Thus we obtain
\begin{align}\label{function_c_m_d_klaus}
\left(\Delta_v \bg \right)_u 
&= 
\sum_{w\subseteq v} (-1)^{|w|} \gamma_{u\cup w} =
\sum_{w\subseteq v} (-1)^{|w|} \gamma_{\varphi(x_{-w}:y_w)}
\nonumber\\
&=
\sum_{w\subseteq v} (-1)^{|w|} f(x_{-w}\!:\!y_w).
\end{align}
Moreover, $x_j = 0 <1 = y_j$ for $j \in v$. 
Due to Lemma~\ref{Lemma32}, $\bg \in \M_d$ 
if $f$ is completely monotonically decreasing, i.e., 
$\M_d^\prime \subseteq \Phi(\M_d)$.

For the proof of the reverse inclusion we extend the mapping
$\varphi$ from $\{0,1\}^d$ to $[0,1]^d$ by
\[
\varphi(x) := \{ j \in [d] : x_j > 0\}.
\]
For $\bg \in \M_d$ let $f \colon [0,1]^d \to {[0,\infty)}$
be defined by 
\[
f(x) := \gamma_{\varphi(x)}
\]
for every $x \in [0,1]^d$. Obviously, we have \eqref{g88}. Thus it
remains to show that $f$ is completely monotonically decreasing.

For $x \in [0,1]^d$ we put 
$\lceil x \rceil := (\lceil x_1 \rceil, \dots, \lceil x_d \rceil) 
\in \{0,1\}^d$. Clearly $\varphi(x) = \varphi(\lceil x \rceil)$,
and therefore $f(x) = f(\lceil x \rceil)$. It follows that $f$ is 
completely monotonically decreasing, if 
\[
\alpha := \sum_{w \subseteq v} (-1)^{|w|} f(x_{-w}\!:\!y_w) \geq 0
\]
for all non-empty $v\in \U_d$ and $x, y\in \{0,1\}^d$ 
such that $x_j \leq y_j$ for every $j \in v$, see Lemma~\ref{Lemma32}.
For $v,x,y$ chosen in this way we put
\[
v_1 := \{j \in v : x_j < y_j\} = \{j \in v : x_j = 0,\ y_j=1\}
\]
and $v_2 := v \setminus v_1$.
Let $w_i \subseteq v_i$ for $1\le i\le 2$. We have
\[
(x_{-(w_1 \cup w_2)}\!:\! y_{w_1 \cup w_2}) = (x_{-w_1}\!:\!y_{w_1}),
\]
which implies
\begin{align*}
\alpha &=
\sum_{w_1 \subseteq v_1} \sum_{w_2 \subseteq v_2}
(-1)^{|w_1|+|w_2|} f(x_{-(w_1\cup w_2)}\!:\!y_{w_1\cup w_2})\\
&=
\sum_{w_1 \subseteq v_1} (-1)^{|w_1|} f(x_{-w_1}\!:\!y_{w_1})
\sum_{w_2 \subseteq v_2} (-1)^{|w_2|}.
\end{align*}
Therefore $\alpha = 0$ if $v_2 \neq \emptyset$, so that
it remains to consider the case $v_2 = \emptyset$, i.e., $v = v_1$.
Without loss of generality we may assume that $y_j = 0$ for
$j \not\in v$, so that $v = \varphi(y)$. Let $u := \varphi(x)$.
Since $u \cap v = \emptyset$, the first part of the proof is
applicable, and \eqref{function_c_m_d_klaus} together with $\bg \in
\M_d$ yields $\alpha \geq 0$. This concludes the proof of
$\Phi(\M_d) \subseteq \M_d^\prime$.
\end{proof}

\end{document}
